\section{Task 4 -- Triển khai Sprint 1}\label{ch:sprint1}

\textbf{Sprint Goal:} có nền tảng kỹ thuật (repository, CI, database có migration và seed) và MVP đầu tiên chạy được cho luồng xem menu.

Sprint 1 diễn ra ngày 15/11/2026. Năm user story được làm (bảng \ref{tab:backlogs1}).

{\footnotesize
\begin{longtable}{@{}L{1.1cm}L{4.3cm}L{1.7cm}L{1.6cm}L{0.9cm}L{1.7cm}L{2.6cm}@{}}
\caption{Product Backlog Sprint 1}\label{tab:backlogs1}\\
\toprule
\textbf{Mã} & \textbf{User story} & \textbf{Yêu cầu} & \textbf{MoSCoW} & \textbf{Point} & \textbf{Người làm} & \textbf{Trạng thái} \\ \midrule \endfirsthead
\toprule \textbf{Mã} & \textbf{User story} & \textbf{Yêu cầu} & \textbf{MoSCoW} & \textbf{Point} & \textbf{Người làm} & \textbf{Trạng thái} \\ \midrule \endhead
\bottomrule \endfoot
US-01 & Là khách hàng, tôi muốn xem danh sách món với giá và tình trạng còn hàng để chọn đồ uống & FR-01, FR-02, NFR-05 & Must & 5 & D (giao diện), A (API) & Done (PR \#11, \#17) \\
US-02 & Là khách hàng, tôi muốn xem chi tiết món với size, topping, mức đường và đá để biết lựa chọn của mình & FR-03, FR-16 & Must & 3 & A & Done -- API (PR \#17); giao diện chi tiết thuộc US-05 \\
US-03 & Là nhóm phát triển, tôi muốn có database với migration và dữ liệu mẫu để chạy và kiểm thử hệ thống giống nhau ở mọi máy & NFR-10 & Must & 3 & B & Done (PR \#15) \\
US-04 & Là nhóm phát triển, tôi muốn CI tự chạy lint, test, build và chặn merge khi lỗi để giữ \code{main} luôn chạy được & NFR-08, NFR-11 & Must & 3 & A & Done (PR \#9, \#10) \\
US-05 & Là nhóm phát triển, tôi muốn có Product Backlog và bảng Issues cho Sprint 1, 2 để theo dõi tiến độ & -- & Must & 2 & C & Done (PR \#14) \\
\end{longtable}
}

\subsection{Task 4.1 -- Thiết lập repository}
Repository \code{tthan-prog/\allowbreak{}BrewLite-AS2} là repository mới, độc lập hoàn toàn với Bài tập lớn 1 (ràng buộc R5). PR \#9 khởi tạo monorepo theo ADR 001:

\noindent\textit{Cấu trúc thư mục chính của repository}
\begin{Verbatim}
apps/
  api/            NestJS API, Prisma schema, migration, seed, Dockerfile
  web/            Next.js (mobile-first), Dockerfile
docs/
  requirements/   Task 1: bối cảnh, SRS, use case
  diagrams/       Task 2: activity, sequence, class, ERD
  architecture/   Task 3: kiến trúc, ADR, deployment diagram
  api/            hợp đồng API
  scrum/          Product Backlog, Sprint Review và Retrospective
  release/        release notes
  deploy/         runbook triển khai
  traceability-matrix.md, deployment.md, contributing.md
tests/e2e/        test E2E Playwright (Sprint 3)
.github/workflows/  ci.yml, cd.yml
docker-compose.yml, .env.example, README.md, CHANGELOG.md
\end{Verbatim}

Các tệp chung ở thư mục gốc:
\begin{itemize}
\item \code{README.md}: hướng dẫn cài đặt, chạy bằng Docker Compose hoặc từng phần, chạy test, build, deploy, rollback.
\item \code{.env.example}: liệt kê biến môi trường với giá trị mẫu, không có giá trị thật.
\item \code{docs/\allowbreak{}contributing.\allowbreak{}md}: chiến lược nhánh, quy ước pull request, commit và Definition of Done (mục \ref{sec:quyuoc}, \ref{sec:dod}).
\end{itemize}

\subsection{Task 4.2 -- Quản lý tài liệu}
Toàn bộ tài liệu nằm trong \code{docs/} và đi qua pull request như mã nguồn. Product Backlog được viết theo user story, có mã yêu cầu, MoSCoW, story point và tiêu chí chấp nhận (PR \#14, Hà Thiên Long đưa lên).

Mỗi user story của Sprint 1 có một GitHub Issue cùng mã, gắn Milestone \code{Sprint 1}. Cuối Sprint, cả 5 Issue đều Closed (hình \ref{fig:issues}).

\anh{t4-16-issues-milestone-sprint1-5-closed.jpg}{0.95\linewidth}{GitHub Issues lọc theo Milestone Sprint 1: 5 Issue US-01 đến US-05 đã đóng\label{fig:issues}}

\subsection{Task 4.3 -- MVP Sprint 1}

\subsubsection{Database, migration và seed (US-03)}
Hồ Thiên Ngân hiện thực Prisma schema theo ERD (PR \#15):
\begin{itemize}
\item 17 model, migration đầu tiên \code{20261115085530\_\allowbreak{}init}.
\item Script \code{prisma/\allowbreak{}seed.\allowbreak{}ts} chạy lặp lại được, tạo: 2 tài khoản, 10 nguyên liệu, 6 sản phẩm, 3 topping, 3 mã khuyến mãi, 2 tài khoản thanh toán thử, 2 phần thưởng.
\end{itemize}

Hình \ref{fig:seed} là kết quả \code{prisma migrate reset} trên PostgreSQL local, chụp trên máy của Ngân.

\anh{t4-08-pr15-prisma-migrate-reset-seed.jpg}{0.85\linewidth}{\code{prisma migrate reset}: áp migration \code{init}, sinh Prisma Client và chạy seed (PR \#15)\label{fig:seed}}

\subsubsection{API sản phẩm (US-01, US-02)}
PR \#17 thêm \code{GET /products} và \code{GET /\allowbreak{}products/\allowbreak{}:id}.

Danh sách món trả \code{maxQuantity}: số ly còn đặt được với lựa chọn mặc định, tính từ công thức và tồn kho nguyên liệu. API không trả số tồn kho thật.

Chi tiết món trả:
\begin{itemize}
\item size kèm phụ phí;
\item topping kèm tình trạng còn hay hết;
\item các mức đường và đá.
\end{itemize}

Món không tồn tại trả 404, id không phải số trả 400. \code{GET /health} trả \code{status: ok} và \code{db: up} sau khi chạy \code{SELECT 1}.

\haianh{t4-11-health-db-up.jpg}{t4-12-get-products.jpg}{\code{GET /health}: API kết nối được PostgreSQL}{\code{GET /products}: 6 món từ seed, có \code{available}, \code{maxQuantity}, không lộ tồn kho}{Kiểm tra API sản phẩm trên máy local}

\subsubsection{Trang menu (US-01)}
Giản Thiên Lực làm trang \code{/menu} (PR \#11):
\begin{itemize}
\item lấy dữ liệu từ API thật;
\item có ba trạng thái: đang tải, danh sách rỗng, lỗi (FR-02);
\item có component test;
\item dùng được ở màn hình 360--414 px (NFR-05).
\end{itemize}

\anh{t4-14-menu-6-mon-tu-api.jpg}{0.95\linewidth}{Trang \code{/menu} ở Sprint 1 hiển thị 6 món lấy từ API}

\subsection{Task 4.4 -- Continuous Integration}
Workflow \code{.\allowbreak{}github/\allowbreak{}workflows/\allowbreak{}ci.\allowbreak{}yml} chạy cho mọi pull request vào \code{main} và mỗi lần push lên \code{main}. Ở Sprint 1, job gồm các bước:
\begin{center}
Checkout $\rightarrow$ Setup Node.js $\rightarrow$ Install (\code{npm ci}) $\rightarrow$ Lint $\rightarrow$ Unit Test $\rightarrow$ Build
\end{center}

Ruleset \code{protect-main} áp lên nhánh \code{main} (4 luật, danh sách bypass trống). Các luật này chặn push thẳng lên \code{main}, buộc mọi thay đổi đi qua pull request, và bắt buộc check ``CI / Install, Lint, Unit Test, Build'' xanh mới được merge.

\anh{t4-03-ruleset-protect-main-active.jpg}{0.9\linewidth}{Ruleset \code{protect-main} đang hoạt động trên nhánh \code{main}}

Đề yêu cầu bằng chứng CI cả thành công lẫn thất bại:
\begin{itemize}
\item \textbf{PR \#9 xanh:} CI chạy hết các bước trong 44 giây (hình \ref{fig:ci9}).
\item \textbf{PR \#10 đỏ:} nhánh \code{test/\allowbreak{}ci-fail-demo} cố ý sửa một assertion trong \code{health.\allowbreak{}controller.\allowbreak{}spec.\allowbreak{}ts}. CI thất bại ở bước Unit Test, và nút Merge bị khóa do check bắt buộc (hình \ref{fig:ci10}, \ref{fig:ci10log}).
\end{itemize}

PR \#10 không được merge, đúng mục đích minh chứng.

\haianh{t4-01-pr9-ci-xanh-all-checks-passed.jpg}{t4-06-pr10-ci-do-merge-bi-khoa.jpg}{PR \#9: CI xanh, ``All checks have passed''}{PR \#10: CI đỏ, nút merge bị khóa}{CI thành công (PR \#9) và thất bại chặn merge (PR \#10)\label{fig:ci9}\label{fig:ci10}}
\anh{t4-07-pr10-log-unit-test-loi.jpg}{0.9\linewidth}{Log của PR \#10: bước Unit Test thất bại tại \code{health.\allowbreak{}controller.\allowbreak{}spec.\allowbreak{}ts:16}\label{fig:ci10log}}

\subsection{Sprint Review và Retrospective Sprint 1}
\textbf{Kết quả:} đạt Sprint Goal, 5/5 user story, 16/16 story point. Demo gồm:
\begin{itemize}
\item API chạy local với PostgreSQL;
\item \code{/products} trả 6 món;
\item trang \code{/menu} lấy dữ liệu thật;
\item CI xanh ở PR \#9 và đỏ ở PR \#10.
\end{itemize}

\textbf{Điều làm tốt:}
\begin{itemize}
\item Làm đúng thứ tự Task 1 $\rightarrow$ 4.
\item Mọi thay đổi qua pull request và CI.
\item Mỗi người commit phần của mình bằng tài khoản riêng.
\end{itemize}

\textbf{Điều chưa tốt:}
\begin{itemize}
\item \textbf{Tạo phẩm phải sửa nhiều lần:}
  \begin{itemize}
  \item ảnh sơ đồ là PNG nhưng đặt đuôi \code{.svg};
  \item \code{erd.puml} lỗi cú pháp;
  \item trang menu dùng class Tailwind khi chưa cài Tailwind.
  \end{itemize}
\item \textbf{Lỗi môi trường:}
  \begin{itemize}
  \item Jest không chạy với NestJS 12, nên chuyển sang Vitest;
  \item Windows Smart App Control chặn binary của SWC, nên bỏ SWC và dùng \code{@Inject} tường minh.
  \end{itemize}
\item \textbf{Quy trình:} nhiều PR được merge khi chưa có approve.
\end{itemize}

\textbf{Hành động cho Sprint 2:}
\begin{itemize}
\item chạy lint, test, build trên máy trước khi push;
\item mở PR nhỏ;
\item có người review trước khi merge;
\item thống nhất hợp đồng API trước khi làm giao diện.
\end{itemize}

\textbf{Phản hồi của Product Owner:} Chấp nhận US-01 đến US-05 theo tiêu chí trong \code{product-backlog.md}. Sprint đạt Sprint Goal: API và trang \code{/menu} chạy được với dữ liệu thật từ database, CI chặn merge khi lỗi. Ghi nhận việc bản backlog đầu chép kế hoạch cũ đã được sửa trước khi chốt. Ưu tiên cho Sprint 2: hoàn thành giỏ hàng, đặt đơn và thanh toán để có luồng đặt món end-to-end.